\begin{proposition}[An explicit locus of algebraic dimension one]\hypertarget{tag-01QX}{}\label{fam:src:D.1.7}\label{fam:st:D.1.7}
\StatementRecord{fam:src:D.1.7}{proved}{01QX}{\SeriesRef{fam:eq:cd-NS}, \SeriesRef{fam:src:D.1.6}, \SeriesRef{fam:src:C.2.1}, \SeriesRef{fam:src:D.1.2}}

There is an explicit nonempty locus of admissible members with ordinary
fibres of algebraic dimension one. On a sufficiently small punctured
cusp neighbourhood set
\[
 c(z)=-6\tau(z)-12\mu(z)-\beta^{\rm part}(z).
\]
Outside a countable set of $z$, the fibre of $X_{c(z)}$ over $\pi(z)$ has
N\'eron--Severi lattice $\mathbb Z\eta+\mathbb ZE$, where the nonzero upper
entries of $E$ are
$E_{13}=-6$, $E_{14}=-6$, $E_{23}=1$, $E_{24}=1$.
Its algebraic dimension is one. The section $P(\pi(z))$ can simultaneously
be required to be non-torsion.
\end{proposition}


\begin{proof}
\begin{enumerate}[label=\arabic*.,leftmargin=*]
\item At the specified parameter $\beta=-6\tau-12\mu$. Both $E$ and $\eta$
satisfy \eqref{fam:eq:cd-NS}, and exact multiplication gives
$\operatorname{Pf}(a\eta+bE)=6a^2$.
For any further integral class $F$, substitution in its period equation
gives a linear relation among $1,\tau,\mu,(\mu+\tau)^2$.
These four functions are linearly independent: applying $g_0$ to a
relation $A+B\tau+C\mu+D(\mu+\tau)^2=0$ and subtracting yields
$-B+D(1-2(\mu+\tau))=0$. A further cusp difference forces $D=0$,
and then $B=0$; nonconstancy of $\mu$ forces $A=C=0$.
Comparison of coefficients now gives
$F_{34}=F_{12}=0$, $F_{13}=-6F_{24}$, and
$6F_{23}-F_{14}=12F_{24}$, exactly
$F\in\mathbb Z\eta+\mathbb ZE$. Every other integral class therefore
has a discrete zero set along this period locus. Excluding their
countable union gives the asserted exact N\'eron--Severi lattice.
It contains the nonzero isotropic $E$ and no matrix of negative
Pfaffian; \ref{fam:src:D.1.6} gives algebraic dimension one.
Finally
$c=-5\tau-12\mu-(\beta^{\rm part}+\tau)$ has imaginary part tending
to minus infinity at the cusp. A sufficiently small cusp neighbourhood
is therefore admissible by \ref{fam:src:C.2.1}. The torsion locus of $P$ is countable
by \ref{fam:src:D.1.2}, so its exclusion still leaves such points in every open subset
of that neighbourhood.
\end{enumerate}
\end{proof}